\documentclass[a4paper]{exam}
\usepackage{amsmath,amssymb,titlesec}
\usepackage[a4paper,margin=22mm]{geometry}

\titleformat{\section}{\normalfont\Large\bfseries}{}{0pt}{}
\titleformat{\subsection}{\normalfont\large\bfseries}{}{0pt}{}
\titlespacing*{\section}{1em}{1.6\baselineskip}{0.6\baselineskip}
\titlespacing*{\subsection}{2em}{1.2\baselineskip}{0.4\baselineskip}

\setlength{\parindent}{0in}
\setlength{\headheight}{18pt}

\printanswers
\title{Week 9}
\author{Radhesh Goel}
\date{\today}

\begin{document}
\maketitle

\section*{Question 1.}

\subsection*{a) Surface area as a function of radius}

\begin{solution}
  The fixed-volume constraint and its rearrangement for height are
  \[
    V=\pi r^2h,
    \qquad
    h=\frac{V}{\pi r^2}.
  \]
  Substituting this into the surface-area equation,
  \[
    \begin{aligned}
      A(r)
      &=2\pi r\left(\frac{V}{\pi r^2}\right)+2\pi r^2 \\
      &=\frac{2V}{r}+2\pi r^2.
    \end{aligned}
  \]
  Therefore, for a physically meaningful radius \(r>0\),
  \[
    \boxed{A(r)=\frac{2V}{r}+2\pi r^2}.
  \]
  For \(V=10\ \mathrm{m^3}\), this becomes
  \(A(r)=20/r+2\pi r^2\).
\end{solution}

\subsection*{b) Bracketing the minimum}

\begin{solution}
  For \(V=10\ \mathrm{m^3}\), begin with \(a=1.5\) and \(b=2.0\). The
  corresponding surface areas are
  \[
    \begin{aligned}
      A(a)&=A(1.5)=\frac{20}{1.5}+2\pi(1.5)^2
            \approx 27.4705\ \mathrm{m^2},\\
      A(b)&=A(2.0)=\frac{20}{2.0}+2\pi(2.0)^2
            \approx 35.1327\ \mathrm{m^2}.
    \end{aligned}
  \]
  Since \(A(b)>A(a)\), the downhill direction is towards decreasing radius.
  Swap \(a\) and \(b\), giving \(a=2.0\) and \(b=1.5\), then calculate
  \[
    c=b+k(b-a)
      =1.5+1.6(1.5-2.0)
      =0.7.
  \]
  At this point,
  \[
    A(c)=A(0.7)=\frac{20}{0.7}+2\pi(0.7)^2
         \approx 31.6502\ \mathrm{m^2}.
  \]
  \par\noindent
  Since \(A(b)<A(a)\) and \(A(b)<A(c)\), the minimum is bracketed. The
  algorithm gives
  \[
    \boxed{a=2.0,\qquad b=1.5,\qquad c=0.7},
  \]
  which corresponds to the interval \(0.7<r<2.0\ \mathrm{m}\), with the
  middle point at \(r=1.5\ \mathrm{m}\).
\end{solution}

\end{document}
